% 核心结论：将A、B移到箭头之外，并明确保持/交换源汇两种对称
% 草图：[真实几何与节点 | 条件、推导、结论]；17cm印宽，最小字号8pt。
% 图契约：端口激励决定对称性如何作用于电势；图中的端口、几何、箭头和方程对应本题。
% 版面：沿用前三模块的几何图+推导分区；目标印宽17cm，简图至中等密度。
% 布局：左侧真实结构，右侧关键关系；真实电路边不作为架构流程箭头处理。
\documentclass[tikz,border=7pt]{standalone}
\usepackage{fontspec,xeCJK,amsmath,amssymb}
\IfFontExistsTF{Times New Roman}{\setmainfont{Times New Roman}}{\setmainfont{TeX Gyre Termes}}
\IfFontExistsTF{Heiti SC}{\setCJKmainfont[AutoFakeBold=1.5]{Heiti SC}}{\setCJKmainfont[AutoFakeBold=1.5]{FandolHei-Regular}}
\usetikzlibrary{arrows.meta,bending,calc,patterns,positioning,decorations.pathmorphing}
\definecolor{ink}{HTML}{243745}
\definecolor{blue}{HTML}{275DAB}
\definecolor{green}{HTML}{177C63}
\definecolor{red}{HTML}{BD4935}
\definecolor{orange}{HTML}{AA6C16}
\definecolor{purple}{HTML}{715393}
\tikzset{every node/.style={font=\fontsize{9.3}{12.5}\selectfont,text=ink},
arrow/.style={-{Stealth[length=5pt 1.5,width'=0pt 3]},shorten >=2pt,shorten <=1pt,line width=1pt},
arrow short/.style={-{Stealth[length=3pt 1.0,width'=0pt 2.5]},shorten >=1pt,shorten <=0pt,line width=.8pt},
axis/.style={arrow,draw=ink!60,line width=.6pt},
vec/.style={arrow,draw=blue},
wire/.style={draw=ink!75,line width=.8pt},
guide/.style={dashed,draw=ink!45,line width=.5pt},
card/.style={rounded corners=2pt,draw=blue!65,fill=blue!4,align=center,text width=6.5cm,inner xsep=9pt,inner ysep=7pt},
title/.style={font=\fontsize{10.5}{14}\selectfont\bfseries,text width=15cm,align=center},
label/.style={font=\fontsize{8.4}{11}\selectfont,align=center},
note/.style={label,text width=6.1cm}}
\begin{document}
\begin{tikzpicture}[x=1cm,y=1cm]
\path[use as bounding box] (-8,-3.5) rectangle(8,3.65);
\node[title] at (0,3.25) {\quad 端口激励决定对称性如何作用于电势};
\begin{scope}[shift={(-4.05,0)}]
\draw[wire,draw=ink](-2,0)--(-1.2305719944826752,0.5001282035862611) (-0.7694280055173248,0.799871796413739)--(0,1.3);\begin{scope}[shift={(-1.0,0.65)},rotate=33.02386755579665]\draw[draw=ink,fill=white,line width=.8pt](-0.275,-.085)rectangle(0.275,.085);\end{scope}
\node[label]at(-1.228895107250074,1.0021463188462676){$r$};
\draw[wire,draw=ink](0,1.3)--(0.7694280055173248,0.799871796413739) (1.2305719944826752,0.5001282035862611)--(2,0);\begin{scope}[shift={(1.0,0.65)},rotate=-33.02386755579665]\draw[draw=ink,fill=white,line width=.8pt](-0.275,-.085)rectangle(0.275,.085);\end{scope}
\node[label]at(1.228895107250074,1.0021463188462676){$r$};
\draw[wire,draw=ink](-2,0)--(-1.2305719944826752,-0.5001282035862611) (-0.7694280055173248,-0.799871796413739)--(0,-1.3);\begin{scope}[shift={(-1.0,-0.65)},rotate=-33.02386755579665]\draw[draw=ink,fill=white,line width=.8pt](-0.275,-.085)rectangle(0.275,.085);\end{scope}
\node[label]at(-0.771104892749926,-0.29785368115373245){$r$};
\draw[wire,draw=ink](0,-1.3)--(0.7694280055173248,-0.799871796413739) (1.2305719944826752,-0.5001282035862611)--(2,0);\begin{scope}[shift={(1.0,-0.65)},rotate=33.02386755579665]\draw[draw=ink,fill=white,line width=.8pt](-0.275,-.085)rectangle(0.275,.085);\end{scope}
\node[label]at(0.771104892749926,-0.29785368115373245){$r$};
\draw[wire,draw=ink](0,1.3)--(0.0,0.275) (0.0,-0.275)--(0,-1.3);\begin{scope}[shift={(0.0,0.0)},rotate=-90.0]\draw[draw=ink,fill=white,line width=.8pt](-0.275,-.085)rectangle(0.275,.085);\end{scope}
\node[label]at(0.42000000000000004,0.0){$r$};
\fill[ink] (-2,0)circle(2pt);\node[below left=4pt] at(-2,0){$A$};
\fill[ink] (2,0)circle(2pt);\node[below right=4pt] at(2,0){$B$};
\draw[vec,green](-2.9,0)--(-2.1,0)node[pos=0,above]{$+I$};\draw[vec,red](2.1,0)--(2.9,0)node[pos=1,above]{$-I$};\node[note]at(0,-2.8){只有保持电阻及端口边界条件的变换才可使用};
\end{scope}
\node[card,text width=6.5cm,draw=blue!65,fill=blue!4]at(3.9,1.8){固定 A 流入、B 流出：$+I,-I$};
\node[card,text width=6.5cm,draw=green!65,fill=green!4]at(3.9,0){保持源汇的镜像：镜像点等势};
\node[card,text width=6.5cm,draw=orange!65,fill=orange!4]at(3.9,-1.8){交换源汇的对称：相对零势电势反号};

\end{tikzpicture}\end{document}